\documentclass{article}
\usepackage{amsmath, amssymb, amsthm}
\usepackage{geometry}
\geometry{a4paper, margin=1in}
\usepackage{fancyhdr}
\usepackage{hyperref}
\usepackage[english, russian]{babel}
\usepackage[utf8]{inputenc}
\usepackage{listings}
\usepackage{xcolor}
\usepackage{tikz}

\pagestyle{fancy}
\fancyhf{}
\fancyhead[L]{Algorithms}
\fancyhead[R]{\thepage}

\title{Algorithms 1 \\ Homework 8}
\author{Nikita Zagainov}
\date{\selectlanguage{english}\today}

\lstset{
    language=Python,
    basicstyle=\ttfamily\small,
    keywordstyle=\color{blue},
    stringstyle=\color{red},
    commentstyle=\color{gray},
    backgroundcolor=\color{lightgray!20},
    frame=single,
    rulecolor=\color{black},
    numbers=left,
    numberstyle=\tiny\color{gray},
    stepnumber=1,
    numbersep=10pt,
    showspaces=false,
    showstringspaces=false,
    breaklines=true,
    breakatwhitespace=true,
    tabsize=4,
    captionpos=b
}

\begin{document}

\maketitle

\begin{center}
	\section*{Обходы и кратчайшие пути в графах}
\end{center}
\begin{enumerate}
	\item Применим алгоритм поиска в глубину к дереву,
	      записывая в каждую вершину время открытия и закрытия.
	      \begin{center}
		      \begin{tikzpicture}
			      \node[draw, circle] (f) at (0,6) {f};
			      \node[draw, circle] (a) at (3,6) {a};
			      \node[draw, circle] (e) at (6,6) {e};
			      \node[draw, circle] (b) at (0,3) {b};
			      \node[draw, circle] (d) at (3,3) {d};
			      \node[draw, circle] (g) at (6,3) {g};
			      \node[draw, circle] (c) at (0,0) {c};
			      \node[draw, circle] (h) at (3,0) {h};

			      \node[above] at (f.north) {8/9};
			      \node[above] at (a.north) {1/10};
			      \node[above] at (e.north) {12/13};
			      \node[above left] at (b.north) {2/5};
			      \node[below] at (d.south) {11/14};
			      \node[above] at (g.north) {3/4};
			      \node[below] at (c.south) {6/7};
			      \node[above] at (h.north) {15/16};

			      \draw[->, >=stealth, line width=1pt] (a) to (f);
			      \draw[->, bend right=15, >=stealth, line width=1pt] (a) to (b);
			      \draw[->, >=stealth, line width=1pt] (a) to (c);
			      \draw[->, >=stealth, line width=1pt] (d) to (c);
			      \draw[->, >=stealth, line width=1pt] (d) to (e);
			      \draw[->, bend left=15, >=stealth, line width=1pt] (f) to (g);
			      \draw[->, bend left=30, >=stealth, line width=1pt] (b) to (g);
		      \end{tikzpicture}
	      \end{center}

	      Далее, отсортируем вершины по убыванию времени закрытия.
	      Результат топологической сортировки:
	      \begin{center}
		      \begin{tabular}{|c|c|c|c|c|c|c|c|}
			      \hline
			      1 & 2 & 3 & 4 & 5 & 6 & 7 & 8 \\
			      \hline
			      h & d & e & a & f & c & b & g \\
			      \hline
		      \end{tabular}
	      \end{center}

	\item Первым шагом, произведём поиск в глубину на
	      инвертированном графе, записывая в каждую вершину время
	      открытия и закрытия:
	      \begin{center}
		      \begin{tikzpicture}
			      \node[draw, circle] (b) at (0,6) {b};
			      \node[draw, circle] (a) at (3,6) {a};
			      \node[draw, circle] (d) at (6,6) {d};
			      \node[draw, circle] (c) at (1.5,3) {c};
			      \node[draw, circle] (e) at (4.5,3) {e};
			      \node[draw, circle] (h) at (0,0) {h};
			      \node[draw, circle] (g) at (3,0) {g};
			      \node[draw, circle] (f) at (6,0) {f};

			      \node[below] at (f.south) {12/15};
			      \node[above] at (a.north) {1/10};
			      \node[above] at (e.north) {4/5};
			      \node[above left] at (b.north) {2/9};
			      \node[above] at (d.north) {11/16};
			      \node[below] at (g.south) {13/14};
			      \node[left] at (c.west) {3/8};
			      \node[below] at (h.south) {6/7};

			      \draw[->, >=stealth, line width=1pt, bend right=10] (a) to (b);
			      \draw[->, >=stealth, line width=1pt, bend right=10] (b) to (a);
			      \draw[->, >=stealth, line width=1pt, bend right=10] (b) to (c);
			      \draw[->, >=stealth, line width=1pt, bend right=10] (c) to (b);
			      \draw[->, >=stealth, line width=1pt] (c) to (a);
			      \draw[->, >=stealth, line width=1pt] (c) to (h);
			      \draw[->, >=stealth, line width=1pt] (g) to (h);
			      \draw[->, >=stealth, line width=1pt] (c) to (e);
			      \draw[->, >=stealth, line width=1pt] (g) to (e);
			      \draw[->, >=stealth, line width=1pt] (f) to (g);
			      \draw[->, >=stealth, line width=1pt, bend right=10] (f) to (d);
			      \draw[->, >=stealth, line width=1pt, bend right=10] (d) to (f);
		      \end{tikzpicture}
	      \end{center}

	      Далее, используя время закрытия вершин, произведём
	      поиск в глубину на исходном графе,
	      строя очередь по убыванию времени закрытия
	      соответствующих вершин в инвертированном графе:

	      \begin{center}
		      \begin{tikzpicture}
			      \node[draw, circle] (b) at (0,6) {b};
			      \node[draw, circle] (a) at (3,6) {a};
			      \node[draw, circle] (d) at (6,6) {d};
			      \node[draw, circle] (c) at (1.5,3) {c};
			      \node[draw, circle] (e) at (4.5,3) {e};
			      \node[draw, circle] (h) at (0,0) {h};
			      \node[draw, circle] (g) at (3,0) {g};
			      \node[draw, circle] (f) at (6,0) {f};

			      \node[below] at (f.south) {2/3};
			      \node[above] at (a.north) {7/12};
			      \node[above] at (e.north) {15/16};
			      \node[above left] at (b.north) {8/11};
			      \node[above] at (d.north) {1/4};
			      \node[below] at (g.south) {5/6};
			      \node[left] at (c.west) {9/10};
			      \node[below] at (h.south) {13/14};

			      \draw[->, >=stealth, line width=1pt, bend left=10] (b) to (a);
			      \draw[->, >=stealth, line width=1pt, bend left=10] (a) to (b);
			      \draw[->, >=stealth, line width=1pt, bend left=10] (c) to (b);
			      \draw[->, >=stealth, line width=1pt, bend left=10] (b) to (c);
			      \draw[->, >=stealth, line width=1pt] (a) to (c);
			      \draw[->, >=stealth, line width=1pt] (h) to (c);
			      \draw[->, >=stealth, line width=1pt] (h) to (g);
			      \draw[->, >=stealth, line width=1pt] (e) to (c);
			      \draw[->, >=stealth, line width=1pt] (e) to (g);
			      \draw[->, >=stealth, line width=1pt] (g) to (f);
			      \draw[->, >=stealth, line width=1pt, bend left=10] (d) to (f);
			      \draw[->, >=stealth, line width=1pt, bend left=10] (f) to (d);
		      \end{tikzpicture}
	      \end{center}

	      Теперь объединить вершины, принадлежащие одному дереву
	      в дереве поиска в компоненты связности.
	      В итоговом конденсате графа будут
	      следующие вершины:
	      \begin{itemize}
		      \item $A: \{a, b, c\}$
		      \item $D: \{d, f\}$
		      \item $E: \{e\}$
		      \item $G: \{g\}$
		      \item $H: \{h\}$
	      \end{itemize}

	      Граф:
	      \begin{center}
		      \begin{tikzpicture}
			      \node[draw, circle] (A) at (1,3) {A};
			      \node[draw, circle] (D) at (6,0) {D};
			      \node[draw, circle] (E) at (4,3) {E};
			      \node[draw, circle] (G) at (3,0) {G};
			      \node[draw, circle] (H) at (0,0) {H};

			      \draw[->] (H) to (A);
			      \draw[->] (H) to (G);
			      \draw[->] (E) to (G);
			      \draw[->] (E) to (A);
			      \draw[->] (G) to (D);
		      \end{tikzpicture}
	      \end{center}

	\item Для определения числа прямых рёбер воспользуемся
	      следующим свойством времени открытия и закрытия вершин:
	      группа из $n$ вершин последовательно соединена $n-1$ прямыми
	      рёбрами, если вершины открыты в существует отрезок времени,
	      в котором лежат все моменты открытия и закрытия вершин, причём
	      время открытия первой и закрытия последней приходятся на одну
	      вершину. Число перекрёстных рёбер равно числу "разрывов" в
	      этой цепи

	      Алгоритм:
	      \begin{lstlisting}
def count_tree_cross_edges(d: list[str | None], f: list[str | None]) -> tuple[int, int]:
	tree_count = 0
	cross_count = -1
	cur_root = None

	for i in range(len(d)):
		if cur_root is None:
			cross_count += 1
			cur_root = d[i]
		elif cur_root == f[i]:
			cur_root = None
		elif d[i] is not None:
			tree_count += 1

	return tree_count, cross_count
		  \end{lstlisting}

	      Сложность - $O(\left| V \right|)$, так как длины массивов
	      $d$ и $f$ равны удвоенному числу вершин в графе.

	\item Алгоритм Дейкстры рассматривает каждую вершину ровно
	      один раз, когда до соответствующей вершины доходит
	      очередь. Это равносильно тому, что у каждой вершины ровно
	      один родитель (у корня - нет родителя). Следовательно,
	      в итоговом графе будет ровно $n-1$ рёбер, а значит,
	      этот граф будет деревом.

	\item \begin{enumerate}
		      \item
		            Предлагаемый алгоритм - алгоритм поиска компонент
		            сильной связности в данном графе:
		            \begin{lstlisting}
def inverse_graph(graph: dict[str, list[str]]) -> dict[str, list[str]]:
	reversed_graph = {}

	for node in graph:
		for neighbor in graph[node]:
			if neighbor not in reversed_graph:
				reversed_graph[neighbor] = list()
			reversed_graph[neighbor].append(node)

	return reversed_graph


def _dfs(
	graph: dict[str, list[str]],
	node: str,
	visited: set[str],
	start: list[str],
	finish: list[str],
):
	visited.add(node)
	start.append(node)
	finish.append(None)
	for neighbor in graph.get(node, []):
		if neighbor not in visited:
			_dfs(graph, neighbor, visited, start, finish)
	start.append(None)
	finish.append(node)


def dfs(graph: dict[str, list[str]]) -> tuple[list[str | None], list[str | None]]:
	visited = set()
	start = list()
	finish = list()

	for node in graph:
		if node in visited:
			continue
		_dfs(graph, node, visited, start, finish)

	return start, finish


def graph_condensates(graph: dict[str, list[str]]) -> list[list[str]]:
	_, _finish = dfs(inverse_graph(graph))
	order = [node for node in reversed(_finish) if node is not None]

	visited = set()
	condensates = list()

	start = list()
	finish = list()

	for node in order:
		if node in visited:
			continue
		_dfs(graph, node, visited, start, finish)
		condensates.append([n for n in finish if n is not None])
		start.clear()
		finish.clear()

	return condensates
	\end{lstlisting}

		            Сложность алгоритма -
		            $O(\left| V \right| + \left| E \right|)$,
		            так как используется инвертирование графа за
		            $O(\left| V \right| + \left| E \right|)$ и
		            поиск в глубину на инвертированном графе за
		            $O(\left| V \right| + \left| E \right|)$.
		            Этот алгоритм корректен, так как он находит
		            компоненты сильной связности в графе по построению.

		      \item Для решения этой задачи, рассмотрим граф из
		            компонент сильной связности, построенный в предыдущем пункте.
		            этот граф не имеет циклов. Обозначим вершины этого графа
		            как истоки, если к ним не ведёт ни одно ребро из других
		            и как стоки, если из них не ведёт ни одно ребро в другие.
		            Тогда минимальное количество рёбер, которое нужно добавить
		            до связности графа будет получено путём соединения
		            стоков и истоков таким образом, чтобы получился один
		            общий цикл. В случае, если в графе есть подграфы,
		            имеющие больше одного стока, эти подграфы можно разбить
		            на графы-пути, начинающиеся с истока и заканчивающиеся
		            в каждом из стоков. Тогда циклично соединив эти графы,
		            мы получим связный граф с минимальным числом добавленных рёбер,
		            равным общему числу стоков.

		            Алгоритм:
		            \begin{lstlisting}
def split_graph(graph: Dict[Any, List[Any]]) -> List[Dict[Any, List[Any]]]:
	components = list()
	visited = set()

	for root in graph:
		if root in visited:
			continue

		stack = [root]

		while stack:
			if stack[-1] not in graph:
				components.append(copy(stack))
			added = False
			for neighbor in graph.get(stack[-1], []):
				if neighbor not in visited:
					stack.append(neighbor)
					added = True
					break
			visited.add(stack[-1])

			if not added:
				stack.pop()

	return components


def add_edges_to_connected(graph: Dict[Any, List[Any]]) -> Dict[Any, List[Any]]:
	condensates = graph_condensates(graph)
	if len(condensates) == 1:
		return graph

	node_to_condensate = dict()
	condensate_to_node = dict()
	for idx, condensate in enumerate(condensates):
		for node in condensate:
			node_to_condensate[node] = idx

		condensate_to_node[idx] = node

	condensed_graph = dict()
	for source, neighbors in graph.items():
		source_condensate = node_to_condensate[source]
		for neighbor in neighbors:
			destination_condensate = node_to_condensate[neighbor]
			if (
				source_condensate != destination_condensate
				and destination_condensate
				not in condensed_graph.get(source_condensate, set())
			):
				if source_condensate not in condensed_graph:
					condensed_graph[source_condensate] = set()
				condensed_graph[source_condensate].add(
					destination_condensate
				)

	condensed_components = split_graph(condensed_graph)

	for i in range(len(condensed_components)):
		source = condensate_to_node[condensed_components[i - 1][-1]]
		destination = condensate_to_node[condensed_components[i][0]]
		if source not in graph:
			graph[source] = list()
		graph[source].append(destination)

	return graph
				\end{lstlisting}

		            Сложность алгоритма - $O(\left| V \right| + \left| E \right|)$.
		            Первым делом он за $O(\left| V \right| + \left| E \right|)$
		            находит конденсаты, далее за
		            $O(\left| V \right| + \left| E \right|)$
		            он строит граф конденсатов, а затем за
		            $O(\left| V \right| + \left| E \right|)$
		            с помощью поиска в глубину он разбивает граф
		            на графы-пути, и, наконец, за
		            $O(\left| V \right|)$ он строит новые рёбра

	      \end{enumerate}

	\item Эту задачу можно свести к обходу в глубину
	      следующим образом:
	      В каждой комнате в комнате на открытии
	      ставим метку времени с помощью специфичного
	      количества монет, соответствующего времени
	      открытия, в случае закрытия --- вторую
	      кучу монет. Если во время обхода мы
	      обнаружили выход, то мы выиграли. Если
	      поиск завершился в исходной комнате, то
	      мы проиграли. Для того, чтобы понимать, из
	      какой комнаты мы пришли в нынешнюю, оставляем
	      монету в соответствующем коридоре.

	      Сложность алгоритма --	-
	      $O(\left| V \right| + \left| E \right|) =
		      O(\left| E \right| + \left| E \right|) =
		      O(\left| E \right|) = O(m)$,

	\item Для начала отсортируем новые рёбра по убыванию
	      координаты истока. Далее, итерируемся по рёбрам,
	      постоянно храня наименьшую координату конца ребра
	      из тех рёбер, которые мы уже рассмотрели.
	      Когда мы рассматриваем очередное ребро, добавляем
	      к количеству компонент расстояние от начала этого
	      ребра до той наименьшей координаты, если начало
	      нового ребра ближе к началу графа-пути, затем
	      обновляем наименьшую координату конца ребра.

	      Алгоритм:
	      \begin{lstlisting}
def get_num_condensates(n: int, new_edges: list[tuple[int, int]]) -> int:
	new_edges.sort(reverse=True)

	n_components = 0
	cur_head = n

	for edge in new_edges:
		n_components += max(cur_head - edge[0], 0)
		cur_head = min(cur_head, edge[1])

	return n_components + cur_head
		  \end{lstlisting}

	      Сложность алгоритма - $O(m \log m)$, так как
	      происходит сортировка рёбер, а затем проход по массиву.

	\item Очевидно, что граф с такими свойствами является
	      множеством цикличных графов, являющихся замкнутыми цепочками
	      из чередующихся вершин и рёбер чётной длины. Для такого графа
	      максимальное паросочетание будет состоять из рёбер,
	      взятых через одно. Такое множество можно получить с
	      помощью поиска в глубину:
	      \begin{lstlisting}
def _dfs(
    graph: dict[str, list[str]],
    node: str,
    matchings: set[tuple[str, str]],
    visited: set[str],
    parent_matched: bool = False,
):
    visited.add(node)

    for neighbor in graph[node]:
        if neighbor in visited:
            continue

        if not parent_matched:
            matchings.add((node, neighbor))

        _dfs(
            graph,
            neighbor,
            matchings,
            visited,
            not parent_matched,
        )


def get_max_matching(graph: dict[str, list[str]]) -> set[tuple[str, str]]:
    matchings = set()
    visited = set()

    for node in graph:
        assert len(graph[node]) == 2
        if node in visited:
            continue

        _dfs(graph, node, matchings, visited)

    return matchings
		  \end{lstlisting}

	      Сложность алгоритма - 
		  $O(\left| V \right| + \left| E \right|)$.
\end{enumerate}

\end{document}


